\chapter{Future Work}

The current study proves a compact audit workflow, not a scalable benchmark
pipeline. It covers one OpenAI API/app pair, a frozen 60-prompt bank, API-only
vendor baselines, API-side repetition checks, and a targeted ChatGPT-app
repetition check for the 17 final cases. The next step is to keep the same
evidence standard while removing the manual app-collection bottleneck.

\section{Highest-Value Extension: Automated App-Side Stability}

The most useful empirical extension is broader app-side repetition with partial
automation. The 17-case targeted stability panel has already been collected; a
follow-up should repeat the full 60-prompt ChatGPT-app collection, then scale to a
pre-registered larger bank. Browser automation with Playwright or Selenium should
open fresh chats, paste prompts, wait for completion, save screenshots, capture
DOM text where available, run OCR fallback, and write an evidence manifest with
visible model labels, tool indicators, and source surfaces.

The output should be a compact prompt-level table: stable, partially stable, or
unstable. Stable means the same divergence label appears again. Partially stable
means the same direction recurs but at least one run changes label. Unstable
means the divergence disappears or reverses. This directly answers the main
reader question beyond the current targeted panel: whether non-divergent rows
remain non-divergent and whether the observed stability pattern holds outside
the final-case subset.

The automation should preserve the current evidence discipline. Each new row
still needs a screenshot, transcript, visible UI-surface note, and final coding
decision. Automation without screenshot QA would add volume without solving the
auditability problem.

\section{Broader Reliability Work}

The completed independent coding pass is sufficient for the present exploratory
thesis, but a larger follow-up should broaden reliability rather than stop at one
blind coder. The next pass should double-code a larger random subset, include
additional app repetitions, and pre-register disagreement-resolution rules before
looking at the results.

The main improvement should target \texttt{safety\_framing\_strength}. That
label had weak agreement in the current study, so future work should either
simplify it into coarser categories or replace it with more objective indicators,
such as explicit refusal, explicit redirection, cited authority, crisis-resource
inclusion, or visible source/tool use. The practical lesson is not that the
thesis becomes invalid, but that safety framing should remain a supporting
descriptor unless the coding guide is tightened and validated.

\section{Prompt Coverage and Power}

A confirmatory version needs larger category cells, a pre-specified primary
endpoint, and a fixed multiple-comparison plan. Scaling from 60 to several
hundred prompts is the direct way to address the current McNemar power problem,
but that scale is practical only if app collection is automated and screenshot
QA is made routine.

The most useful expansion would add objective instruction-following prompts and
more balanced refusal-boundary subcategories. Objective format constraints such
as JSON validity, exact separators, forbidden punctuation, and token counts are
valuable because they are easier to audit than soft safety framing. Refusal
prompts should remain redacted in shared artifacts and balanced across risk
families so that one domain does not drive the result.

\section{Other Providers and Longitudinal Drift}

Claude and Gemini are currently API-only baselines. They contextualize the
OpenAI result but do not support Claude-app or Gemini-app claims. A direct
extension would collect official consumer-app outputs for those providers using
the same screenshot protocol and coding guide. Claims should then be framed as
provider-specific API/app pairs, not as a single leaderboard.

The study should also be repeated over time. LLM products can change model
routing, search behavior, visible citations, memory, and safety wording without
obvious public notice. A longitudinal audit could rerun the same prompt bank
weekly or monthly and report prompt-level changes, not only aggregate rates. The
important question would be whether the same refusal and factual-support cases
persist across product updates.

\section{Mechanism-Probing Without Causal Overclaiming}

This thesis cannot attribute differences to hidden system prompts, tools,
retrieval, moderation layers, or routing. Future work can move closer to
mechanism only by adding controlled surface ablations. For the app, this means
documenting whether search, memory, location permissions, and visible tool
indicators are enabled. For the API, it means varying system prompts, tool
availability, temperature, and response-format settings. Such ablations can show
which visible conditions are sufficient to reproduce a pattern, but they still do
not reveal proprietary internals.

\section{Priority Order}

\begin{longtable}{L{0.08\linewidth}L{0.26\linewidth}L{0.28\linewidth}L{0.24\linewidth}}
\caption[Prioritized next steps for the deployment-surface audit]{Prioritized next steps for strengthening the deployment-surface audit.}\\
\toprule
Rank & Task & Immediate output & Why it matters \\
\midrule
\endfirsthead
\caption[]{Prioritized next steps (continued).}\\
\toprule
Rank & Task & Immediate output & Why it matters \\
\midrule
\endhead
\bottomrule
\endfoot
1 & Automated full-bank app repetitions & Stability table with screenshots, transcripts, and manifests & Tests whether the targeted stability pattern holds beyond the 17 final cases. \\
2 & Larger automated prompt bank & Better-powered confirmatory audit & Separates durable patterns from small-cell effects. \\
3 & Other consumer apps & Provider-specific API/app matrices & Tests whether the access-surface issue generalizes. \\
4 & Longitudinal waves & Drift report by prompt and category & Measures persistence over product updates. \\
5 & Broader reliability coding & Multi-coder reliability report on a larger subset & Quantifies where the taxonomy remains dependable at scale. \\
6 & Surface ablations & Search, memory, locale, tool, and system-prompt condition notes & Narrows plausible explanations without claiming hidden mechanisms. \\
\end{longtable}

This agenda turns the thesis into a larger engineering program: automate the app
surface, preserve screenshot-level evidence, expand the prompt bank, and then
test stability, reliability, and generality before making stronger claims.
